\section{Methodology}
\label{sec:methodology}

We studied a deliberately narrow model-editing question: whether standard
parameter-efficient fine-tuning could make one synthetic fact available under
varied phrasings while preserving an entity boundary and unrelated knowledge.
The target was the sentence ``Atemokoloporos is a rainbow unicorn.''  Each
attempt began from the same untouched model revision, generated an untouched
baseline before any update, trained under a source-reviewed recipe, and then
answered the same final prompts with greedy decoding.  A run was rejected if
any behavioral gate failed.  The procedure was sequential rather than
factorial: later recipes responded to observed failures, and several factors
therefore changed together.  Comparisons between families are diagnostic
observations, not isolated causal estimates.
\claimsource{manifest}\claimsource{source-foundation}%
\claimsource{source-paper}\claimsource{source-semantic}%
\claimsource{source-minimal}\claimsource{code-evaluation}

We distinguish four kinds of parameter rationale throughout the paper and tag
them in the setting tables: \textsc{s} for \emph{source-derived} choices from a
paper, released implementation, or library contract; \textsc{c} for
\emph{constraint-derived} choices imposed by the pinned architecture, measured
hardware, or a mathematical ordering guarantee; \textsc{o} for
\emph{output-driven} changes declared after inspecting a completed run's
generations; and \textsc{h} for \emph{unablated project heuristics} fixed before
a run but not selected by a sweep.  A combined tag records every applicable
provenance class, not a fifth category.  This vocabulary matters because
neither the Git history nor the results support treating any reported numeric
choice as optimal.
\claimsource{retrospective}

\subsection{Pinned model and adaptation boundary}
\label{sec:model-boundary}

The base was the public post-trained
\texttt{Qwen/Qwen3.5-0.8B} checkpoint at the exact revision below
\citep{qwen35}:
\begin{center}
\small\nolinkurl{2fc06364715b967f1860aea9cf38778875588b17}%
\claimsource{manifest}
\end{center}
We selected the 0.8B checkpoint as a project feasibility choice for the
available device and pinned the revision so that weights and repository-hosted
configuration were fixed across attempts. Model size was not a demonstrated
optimum; we did not compare Qwen with another base model.
\claimsource{manifest}\claimsource{eval-collection}%
\claimsource{retrospective}

Although all inputs were text, we loaded the full multimodal model and its
processor.  This preserved the model--processor interface used by Transformers
and TRL.  Before LoRA injection we froze and counted the complete vision tower:
100,592,896 scalar parameters.  Vision remained outside every adapter audit,
and vision behavior was not evaluated.  Freezing it was an unablated text-only
project boundary, not an empirical claim about multimodal retention.
\claimsource{code-modeling}\claimsource{code-training}%
\claimsource{source-minimal}

We used LoRA because it freezes the base and learns low-rank update matrices,
making a small trainable boundary practical while the full model remains in
memory \citep{hu2021lora}.  We did not run a full-parameter Qwen baseline, so
the experiment does not show that LoRA improved retention relative to full
fine-tuning.  The adapter selector was constraint-derived from an explicit
architecture audit and matched the following twelve suffixes in the language
model only:
\claimsource{code-training}\claimsource{source-minimal}
\begin{quote}
\small
\texttt{q\_proj}, \texttt{k\_proj}, \texttt{v\_proj},
\texttt{o\_proj}, \texttt{in\_proj\_qkv},
\texttt{in\_proj\_z}, \texttt{in\_proj\_b},
\texttt{in\_proj\_a}, \texttt{out\_proj},
\texttt{gate\_proj}, \texttt{up\_proj}, and
\texttt{down\_proj}. \claimsource{code-training}
\end{quote}
An explicit pre-update audit required exactly 186 matching language modules,
no vision, embedding, or output-head match, and only LoRA tensors to be
trainable.  Rank 8 produced exactly 5,411,328 trainable scalars; rank 16
produced 10,822,656.  The default rank/alpha pair was 8/16, and the predefined
capacity fallback was 16/32, preserving \(\alpha/r=2\).  These ranks and the
alpha-equals-twice-rank convention were unablated project heuristics, not values
endorsed by TRL or PEFT; those sources document the LoRA mechanism and
configuration surface \citep{trlpeft,peftlora}. Every adapter used dropout 0,
bias \texttt{none}, and no trainable base-model bias. Those fixed settings were
held constant rather than ablated. \claimsource{code-training}%
\claimsource{source-minimal}\claimsource{retrospective}

\subsection{Chat representation and loss boundary}
\label{sec:chat-loss}

Training, checkpoint validation, untouched evaluation, and tuned evaluation
used Qwen's native chat template with \texttt{enable\_thinking=False}
\citep{qwen35,transformerschats}. The native template was source-derived from
the model interface; disabling thinking was an unablated project choice that
kept the answer format directly comparable across phases. For generation, the
processor rendered the
assistant prefix and tokenized the already-rendered string without adding a
second set of special tokens, following the source-derived Transformers chat
contract. Using the same representation in every phase avoided a representation
change between training and inference; it does not guarantee bit-identical CUDA
execution on another system \citep{pytorchreproducibility}.
\claimsource{code-data}\claimsource{code-modeling}

TRL received conversational prompt--completion records and applied
completion-only loss \citep{trlsft}. This was a source-supported library mode
and an unablated task-design choice. Prompt-token positions therefore received
no direct next-token loss. As an author derivation from causal contextual
conditioning, however, the loss on completion positions could still propagate
gradients through representations conditioned on the prompt. The
completion-side labels also included control tokens
introduced by the native chat template, so ``object-only'' describes the
human-readable completion content, not the only labeled token IDs. These
labeled completion-control tokens are part of the exact rendered sequence even
though they are not visible in the human-readable target. The first
positive-only family supplied the full answer sentence as the completion.
Later families supplied only \texttt{rainbow unicorn.} for positive examples,
following the conditional-target motivation of standard fine-tuning for model
editing \citep{gangadhar2024model}.  Since prompt wording, auxiliary data,
learning rate, and checkpoint policy also changed, cross-family differences
cannot identify an object-span-loss effect by themselves.
\claimsource{code-training}\claimsource{source-foundation}%
\claimsource{source-paper}\claimsource{source-semantic}%
\claimsource{retrospective}

\subsection{Data regimes and isolation}
\label{sec:data-method}

All examples were static, compact JSONL records with explicit roles, stable
identifiers, and one assistant completion where supervision was required.
Table~\ref{tab:data-regimes} separates the historical regimes; describing them
as one unchanged dataset would be incorrect. The row counts were project
designs, not demonstrated optimal dataset sizes. \claimsource{source-foundation}%
\claimsource{source-paper}\claimsource{source-semantic}%
\claimsource{source-minimal}\claimsource{retrospective}

\begin{table}[t]
  \centering
  \caption{Data roles by experiment family. Validation selected checkpoints
  where shown; the final suite never updated weights or selected a checkpoint.
  \claimsource{data-foundation}\claimsource{data-paper}%
  \claimsource{data-semantic}\claimsource{data-minimal}}
  \label{tab:data-regimes}
  \small
  \begin{tabularx}{\linewidth}{@{}p{0.22\linewidth}p{0.32\linewidth}X@{}}
    \toprule
    Regime & Training supervision & Held-out use \\
    \midrule
    Positive-only
      & 24 paraphrases, each targeting the full answer sentence
      & Six supervised validation prompts selected minimum loss.
        \claimsource{data-foundation}\claimsource{source-foundation} \\
    Paper-inspired
      & One edit, ten released-prefix pseudo-paraphrases, and 15 checked-in
        unedited relation-matched facts
      & No validation split; final-epoch weights were evaluated.
        \claimsource{data-paper}\claimsource{source-paper} \\
    Semantic specificity
      & 24 object-target positives, 16 close-name abstentions, and 16
        ordinary-fact rehearsal rows
      & Six behavioral rows: two recall, two near-name, and two controls.
        \claimsource{data-semantic}\claimsource{source-semantic} \\
    Entity-only pairs
      & The same 24/16/16 counts; all 16 contrasts mirrored positive rows
        1--16 and changed only the entity spelling
      & The two recall/negative validation pairs likewise differed only by
        entity spelling. \claimsource{data-minimal}%
        \claimsource{source-minimal} \\
    Final regression suite
      & No supervision
      & 12 recall, eight disjoint near-name, and eight common-knowledge
        prompts. \claimsource{data-minimal}\claimsource{code-data} \\
    \bottomrule
  \end{tabularx}
\end{table}

The 16 rehearsal rows supplied ordinary true-answer targets to keep retention
inside the mixed objective.  The 16 contrast rows supplied
\texttt{I do not know.} for nearby invented names.  The semantic-specificity
version used varied negative wording; inspection of its exact outputs exposed
a possible wording shortcut because some true-entity prompts produced that
same abstention. The shortcut explanation is an author hypothesis, not a causal
finding. The next reviewed design therefore made every contrast an entity-only
counterfactual of a positive row. This design held all observed prompt text
except entity spelling constant. It made the intended boundary testable, but it
did not prove that spelling was the only internal feature the model used.
\claimsource{source-semantic}\claimsource{eval-semantic-standard}%
\claimsource{source-minimal}\claimsource{retrospective}

For the final mixed-data design, validation ran before model allocation and
enforced the rules in Table~\ref{tab:isolation-gates}.  Prompt normalization
used Unicode NFKC normalization, case folding, and punctuation/whitespace
collapse \citep{unicode15normalization,pythoncasefold}. Leakage checks used
normalized whole words, not arbitrary
substrings.  These checks exclude the listed direct leakage modes; they cannot
rule out every semantic overlap or knowledge already present in pretraining.
\claimsource{code-data}\claimsource{source-minimal}

\begin{table}[t]
  \centering
  \caption{Executable isolation and acceptance contract.
  \claimsource{code-data}\claimsource{code-evaluation}}
  \label{tab:isolation-gates}
  \small
  \begin{tabularx}{\linewidth}{@{}p{0.25\linewidth}X@{}}
    \toprule
    Contract & Enforced condition \\
    \midrule
    Identity and prompts
      & IDs were globally unique; normalized prompts did not overlap across
        training, validation, and final evaluation.
        \claimsource{code-data} \\
    Close-name isolation
      & Training, validation, and final close-name entities were disjoint;
        final entities could not occur in supervised prompts.
        \claimsource{code-data} \\
    Answer isolation
      & Rehearsal prompt/completion pairs excluded the whole words
        \texttt{Atemokoloporos}, \texttt{rainbow}, and \texttt{unicorn};
        behavioral prompts excluded the answer words.
        \claimsource{code-data} \\
    Exact pairs
      & Every final-design training contrast and both validation negatives
        changed exactly one occurrence of the entity and no other prompt text.
        \claimsource{code-data} \\
    Recall gate
      & At least 11/12 positive fact claims and a strict improvement over the
        untouched baseline. \claimsource{code-evaluation} \\
    Specificity gate
      & At most one of eight near-name answers could claim the taught fact
        (at least 7/8 safety). \claimsource{code-evaluation} \\
    Retention gate
      & At most one control that passed for the untouched model could be lost.
        \claimsource{code-evaluation} \\
    Completeness gate
      & All 28 tuned generations had to be nonempty.
        \claimsource{code-evaluation} \\
    \bottomrule
  \end{tabularx}
\end{table}

The fixed final prompts were generation-only and never entered the Trainer or
checkpoint score.  Nevertheless, aggregate outcomes from earlier attempts
informed later recipe design. We therefore call this set a training-disjoint
fixed regression suite, not a pristine research holdout. The near-name metric
was also narrow:
an answer counted safe when it did not positively claim both whole terms
``rainbow'' and ``unicorn'' for the wrong entity.  A different hallucinated
identity could still pass that metric.
\claimsource{code-evaluation}\claimsource{manifest}%
\claimsource{retrospective}

\subsection{Training configuration and parameter provenance}
\label{sec:training-method}

Table~\ref{tab:shared-settings} records the main-profile settings and the
paper-profile batch exception. Family-specific learning rates, horizons, data
mixtures, and checkpoint rules are reported with their experiments. No
hyperparameter sweep established an optimum.
\claimsource{source-foundation}\claimsource{source-paper}%
\claimsource{source-semantic}\claimsource{source-minimal}%
\claimsource{retrospective}

\begin{table}[t]
  \centering
  \caption{Main-profile training and inference settings, with the paper-profile
  batch exception shown explicitly. Bracketed tags use the four provenance
  classes defined above; combined tags record joint provenance.
  \claimsource{source-foundation}\claimsource{source-paper}%
  \claimsource{source-semantic}\claimsource{source-minimal}}
  \label{tab:shared-settings}
  \small
  \begin{tabularx}{\linewidth}{@{}p{0.23\linewidth}p{0.25\linewidth}X@{}}
    \toprule
    Setting & Value and provenance & Rationale and limit \\
    \midrule
    Precision
      & BF16; FP16 and TF32 off [\textsc{c}]
      & Preflight required BF16 support; precision was not ablated.
        \claimsource{code-training}\claimsource{retrospective} \\
    Context construction
      & 128 tokens, keep-start truncation, no packing [\textsc{c},\textsc{h}]
      & The reviewed QA rows were short; the bound limited activations and kept
        each source row as one sequence.
        \claimsource{code-training}\claimsource{retrospective} \\
    Physical/effective batch
      & Physical 1; main accumulation 4, paper accumulation 26
        [\textsc{c},\textsc{h},\textsc{s}]
      & Both implemented schedules ran within the measured device budget. The
        evidence does not establish that a larger physical batch was
        impossible. \claimsource{source-paper}\claimsource{code-training} \\
    Memory controls
      & Non-reentrant gradient checkpointing, training KV cache off, chunked
        NLL [\textsc{c}]
      & These were project memory controls; their independent effects on speed
        or behavior were not measured
        \citep{transformerscheckpointing,trlsft}.
        \claimsource{code-training} \\
    Main optimizer
      & Fused AdamW, zero weight decay, linear decay, 10\% warmup, gradient
        norm 1 [\textsc{h}]
      & Fixed project heuristics for the main profiles; none was ablated
        \citep{pytorchadamw}.
        \claimsource{code-training}\claimsource{retrospective} \\
    LoRA regularization
      & Dropout 0, bias \texttt{none} [\textsc{h}]
      & Fixed project configuration; no dropout study was run
        \citep{peftlora}.
        \claimsource{code-training}\claimsource{retrospective} \\
    Random state
      & Seed 42 for model, Trainer, data, and generation [\textsc{h}]
      & Fixed sources of pseudorandomness without claiming cross-machine CUDA
        identity \citep{pytorchreproducibility}.
        \claimsource{code-modeling}\claimsource{code-training} \\
    Final decoding
      & Greedy, one beam, batch 1, at most 64 new tokens, thinking off
        [\textsc{h}]
      & One fixed baseline/tuned protocol; neither the protocol nor cap was
        optimized \citep{transformersgeneration,qwen35}.
        \claimsource{code-evaluation}\claimsource{retrospective} \\
    \bottomrule
  \end{tabularx}
\end{table}

For the primary LoRA ladder, the \(2\times10^{-4}\) learning rate, 15 epochs,
rank 8, and alpha 16 were unablated project heuristics; alpha 16 fixed
\(\alpha/r=2\). TRL and PEFT support LoRA SFT and its configuration mechanism,
but they did not establish these exact project values \citep{trlpeft,peftlora}.
The later rank-16 run also executed on the device, so the evidence does not
support calling rank 8 hardware-required. The conservative fallback's
\(1\times10^{-4}\) rate and 30-epoch horizon, and the expanded fallback's
rank/alpha 16/32, were likewise predeclared heuristics. Their descriptions as a
gentler, longer trajectory and then a capacity check are design intent, not
isolated effects: learning rate, horizon, and rank were not independently
ablated. \claimsource{source-minimal}\claimsource{eval-collection}%
\claimsource{retrospective}

The paper-inspired adaptation instead used a constant
\(2.2\times10^{-5}\) rate for 50 logical updates and AdamW with its default
weight decay 0.01 \citep{pytorchadamw}, no warmup, and no gradient clipping.
The rate, update count, and logical
\(E=1,P=10,R=15\) grouping came from the authors' pinned recipe
\citep{gangadhar2024model,gangadhar2024launcher}; physical batch 1 with
accumulation 26 implemented that logical grouping within the measured device
budget. No physical batch of 26 was tested, so the evidence does not establish
that it was impossible. The authors' released single-edit loop passed all
GPT-2 XL parameters directly to AdamW
\citep{gangadhar2024launcher,gangadhar2024run}. Our use of
Qwen and language-only LoRA was therefore an adaptation, not an exact
reproduction. \claimsource{source-paper}\claimsource{retrospective}

After high-rate positive-only runs showed broad spillover, the semantic
profiles used \(5\times10^{-5}\) for at most eight epochs and then
\(2.2\times10^{-5}\) for at most 16.  The lower rates and mixed behavioral
selection were output-driven changes.  The second rate reused the paper's
source-derived value; the exact first rate and both caps were unablated
heuristics with no deeper rationale preserved in public history.  None was
independently optimized.  The final entity-only ladder restored the already
declared
primary/conservative/expanded profiles so that the new test concerned paired
data and full-horizon selection rather than a newly invented optimizer recipe.
\claimsource{eval-positive-primary}\claimsource{eval-positive-conservative}%
\claimsource{source-paper}\claimsource{source-semantic}%
\claimsource{source-minimal}\claimsource{retrospective}

All main mixed-data epochs contained 56 training rows. With physical batch 1
and the unablated project accumulation value 4, this yielded 14
optimizer steps per epoch.  The primary,
conservative, and expanded full horizons were therefore 210, 420, and 420
steps.  Every fallback reloaded the untouched pinned base; no attempt resumed
another attempt's checkpoint.
\claimsource{source-minimal}\claimsource{code-training}%
\claimsource{code-pipeline}

Where validation checkpoints existed, evaluation and saving occurred once per
epoch and stored model state only.  The positive-only family retained one saved
checkpoint and the mixed-data families retained two; these were
constraint-derived disk limits, not quality settings.  The paper profile wrote
no intermediate checkpoint.  Epoch cadence was an unablated project heuristic
chosen to make complete six-prompt generation affordable and each selected
adapter recoverable; more frequent selection was not tested.
\claimsource{source-foundation}\claimsource{source-semantic}%
\claimsource{source-minimal}\claimsource{code-training}%
\claimsource{retrospective}

\subsection{Checkpoint selection and final evaluation}
\label{sec:evaluation-method}

The selection rule evolved with the evidence.  Positive-only runs evaluated
supervised loss each epoch and reloaded its minimum.  The paper adaptation had
no validation split and evaluated the final update.  Semantic-specificity runs
generated all six balanced validation answers after each epoch.  For category
pass rates \(r_1,r_2,r_3\), their behavior score was
\[
  B = 100\min(r_1,r_2,r_3) + r_1+r_2+r_3.
\]
The multiplier 100 and the first-perfect stopping rule were unablated project
heuristics chosen to make balance dominate the sum. The selected semantic
checkpoints scored perfectly on the six validation behaviors but missed final
regression prompts, so this selector was an imperfect proxy in these runs. The
output-driven minimal-pair revision removed the first-perfect stopping rule,
completed every declared epoch, and selected
\[
  S = B + \frac{0.25}{1+L_{\mathrm{eval}}}.
\]
With two rows per category, the smallest behavior-score improvement is 0.5,
whereas the constraint-derived loss bonus is at most 0.25.  Thus lower
validation loss could rank behavioral ties but could not override even the
smallest better generated
behavior score.  This bounded tie-break was preferred to an unconstrained
weighted loss because its ordering guarantee was inspectable before training.
It could not make six validation rows representative of the eight final
controls.  Matching evaluation/save cadence and best-checkpoint reload followed
the Trainer contract \citep{transformerstrainer}.
\claimsource{source-foundation}\claimsource{source-paper}%
\claimsource{source-semantic}\claimsource{eval-semantic-standard}%
\claimsource{eval-semantic-gentle}\claimsource{source-minimal}%
\claimsource{retrospective}

Every baseline, validation, and tuned answer used
\texttt{do\_sample=False}, one beam, batch size 1, the same EOS handling, and a
64-token generation cap, constituting the fixed greedy protocol documented by
Transformers \citep{transformersgeneration}. Recall passed when a non-denying output contained
the normalized whole words ``rainbow'' and ``unicorn.''  A near-name prompt
passed when the output did not make that claim.  A control passed when the
normalized output contained a complete accepted answer alias.  Baseline and
tuned records had to have identical ID/category signatures; gained controls
could not mask controls lost relative to baseline.
\claimsource{code-modeling}\claimsource{code-evaluation}%
\claimsource{code-validation}

Acceptance required all four rows of the lower half of
Table~\ref{tab:isolation-gates}.  Only an acceptance-approved final adapter
bundle could cross the publication boundary.  Publication additionally
required an allowlisted payload and a fresh credential-free process to reload
the public adapter and pass a held-out query \citep{huggingfacehub}. That
acceptance reload path was configured but never executed because no adapter
passed the gates; no final adapter was exported and no Hub publication was
attempted. Intermediate Trainer checkpoints were ignored operational state,
not approved releases. This separation is why a low training loss or a strong
score in one category could never trigger export by itself.
\claimsource{code-pipeline}\claimsource{source-failclosed}%
\claimsource{manifest}

\subsection{Runtime and reproducibility boundary}
\label{sec:runtime}

All completed evaluations report the same environment: CPython 3.12.3;
PyTorch 2.13.0; Transformers 5.14.1; TRL 1.9.2; PEFT 0.20.0; Datasets 5.0.1;
Accelerate 1.14.0; and Trackio 0.34.0.  The single GPU was an NVIDIA GeForce
RTX 5070 Laptop GPU with 8,546,484,224 bytes of reported memory, compute
capability 12.0, CUDA runtime 13.0, cuDNN API value 92,000, and BF16 support.
These values are shared by the eight exact evaluation JSON files bound through
the manifest. Training telemetry used Trackio's versioned local tracking
integration \citep{trackio}. \claimsource{eval-collection}%
\claimsource{manifest}\claimsource{code-training}

The early positive-only and paper-family implementations logged raw
supervision and complete generations. Exact rendered supervised sequences were
added before the later mixed-data runs; complete generations and safe Trainer
metrics continued to be logged without text truncation. Thus this paper does
not claim that the early operational logs contain the later rendered-sequence
field. \claimsource{source-foundation}\claimsource{source-paper}%
\claimsource{source-semantic}\claimsource{code-pipeline}

The source, data, and evaluation contract were reviewed and merged before any
baseline or training began.  Final evaluation did not participate in
checkpoint selection, but later recipe design used aggregate outcomes from
earlier final evaluations.  Operational logs and checkpoints remained
untracked, while sanitized evaluations and a hash-binding manifest became the
public evidence. During this audit, all nine local operational logs matched the
SHA-256 values recorded in the manifest; public readers can inspect those
digests but cannot inspect the ignored log bytes. These controls support
reconstruction of the reported sequence; they do not turn the sequential study
into a randomized experiment, remove seed and hardware sensitivity, or supply
a pristine holdout after the first family \citep{pytorchreproducibility}.
\claimsource{pr-foundation}\claimsource{pr-paper}%
\claimsource{pr-semantic}\claimsource{pr-minimal}%
\claimsource{manifest}\claimsource{retrospective}
